\chapter{پذیرش دانشگاه}\label{ch:admission}

\section{انتخاب برنامه و دانشگاه}
اولین و شاید مهمترین عاملی که باید در انتخاب رشته و دانشگاه آینده خود در نظر بگیرید علاقه، تخصص و آینده‌ایست که برای خود متصور هستید، 
این کار به خودی خود ساده نیست ولی برای هر شخص متفاوت است.

انتخاب از بین هزاران رشته، دانشگاه و حتی مقطع‌هایی که می‌توان در آن‌ها درس خواند بوسیله معیار‌های مختلف انجام می‌شود. 
این معیار‌ها گاهی برای هر شخص متفاوت هستند، مثلا شاید کشور الف را به کشور ب ترجیح بدهید چون دوستان یا آشناهایی در آن کشور دارید. ولی عواملی 
هم وجود دارند که برای عموم افراد مهم هستند. در این بخش بیشتر به این موارد می‌پردازیم. 

همانطور که قبلا گفته شد با توجه به شرابط نامناسب ویزای کشور‌های مختلف برای دانشجویان ایرانی بهتر است که روی یک کشور و رشته خاص تمرکز نکنید و همه 
تخم‌مرغ‌های خود را در یک سبد نگذارید. 

\subsection{معیارهای انتخاب کشور}
یکی از کلیدی‌ترین تصمیمات در فرآیند مهاجرت تحصیلی، انتخاب کشور مقصد است. کشوری که دانشگاه شما در آن قرار دارد،
 تاثیر مستقیمی بر آینده حرفه‌ای و شخصی شما خواهد داشت. در ادامه، معیارهای اصلی برای انتخاب کشور مقصد آورده شده است:

\subsubsection{آینده شغلی و قوانین ویزا}
کشور مقصد همان جایی است که احتمالا پس از فارغ‌التحصیلی باید در آن \textbf{دنبال شغل بگردید} و \textbf{از آن ویزای کار بگیرید}.
قوانین کار پس از تحصیل و نرخ بیکاری در هر کشور متفاوت است و تاثیر مستقیمی بر امکان ماندن شما پس از تحصیل دارد.
به عنوان دانشجوی ایرانی اگر دارای پاسپورت دوم نیستید، ویزا گرفتن برای تمام کشورها ممکن است کار ساده‌ای نباشد. بنابرین 
توجه به سابقه صدور ویزای کشور مورد نظر شما برای دانشجوهای ایرانی حائز اهمیت است. 

\subsubsection{هزینه‌های زندگی و تحصیل}
شما باید بتوانید \textbf{هزینه‌های زندگی در آن کشور را بپردازید}. این هزینه‌ها
شامل اجاره مسکن، خوراک، حمل و نقل و بیمه در کنار شهریه دانشگاه (در صورت عدم دریافت فاند کامل) 
است که در کشورهای مختلف تفاوت چشمگیری دارد. بنابرین باید هنگام انتخاب کشور به فرصت‌های تحصیلی در آن کشور، فرصت‌های بورسیه و 
توانایی مالی خود توجه کنید. 


\subsubsection{بورسیه‌ها و کمک‌های مالی}
در برخی کشور‌ها، دولت برای تشویق دانشجوهای خارجی به تحصیل در آن کشور بورسیه‌هایی تعیین می‌کنند، این بورسیه‌ها ممکن است 
تمام یا بخشی از هزینه تحصیل را پوشش دهند. هر کشور در این زمینه قوانین خاص خود را دارد. یک مثال مناسب از این بورسیه‌ها بورسیه استانی در 
ایتالیا است که توسط هر استان در ایتالیا به دانشجویان دارای توانایی مالی پایین تعلق می‌گیرد. توجه کنید که این بورسیه‌ها جدا از بورسیه‌های ارائه  
شده توسط دانشگاه‌ها و اساتید هستند. 

\subsubsection{منابع کسب اطلاعات و تحقیق}
برای کسب اطلاعات در این رابطه، گروه‌های مهاجرتی کشور‌های مختلف 
(مثلا گروه‌های تلگرامی پرگار\footnote{\href{https://t.me/pargarGlinks}{لینک گروه‌های پرگار}}) 
و سوال کردن از سال‌بالایی‌ها و دانشجویان مقیم آن کشور می‌توانند منابع خوبی باشند.
ولی به آن‌ها اکتفا نکنید و حتماً تحقیقات مستقل خود را از طریق سایت‌های رسمی دولتی و دانشگاهی انجام دهید.



\subsection{معیارهای انتخاب دانشگاه}

\subsubsection{رتبه و شهرت}
تحصیل در دانشگاه‌های بهتر به شما فرصت‌های تحقیقاتی و کاری بهتری می‌دهد. فارغ از نوع مسیری که برای آینده خود انتخاب می‌کنید که می تواند کاری یا آکادمیک باشد 
تحصیل در یک دانشگاه بهتر یا دانشگاهی که در رشته شما شهرت بهتری دارد با بهبود سابقه و رزومه شما در آینده به شما کمک خواهد کرد\footnote{درباره ویزای \lr{UK High Potential Individual} سرچ کنید}. واضح است که دانشگاه تنها
عامل تعیین کننده فرصت‌های شغلی و آکادمیک‌ شما در آینده نیست و کشور، استاد و فعالیت‌های خودتان در مقطعی که در آن درس میخوانید میتواند حتی مهمتر باشد. 
با این حال منابع زیر می‌تواند برای بدست آوردن اطلاعات درباره رشته‌های مختلف مفید باشد:

\begin{itemize}[leftmargin=1.5cm]
  \item \href{https://www.topuniversities.com/university-rankings}{\lr{QS Ranking}}
  \item \href{https://www.timeshighereducation.com/world-university-rankings}{\lr{Times Higher Education Ranking}} 
  \item \href{https://www.shanghairanking.com/}{\lr{Shanghai Ranking}}
  \item \href{https://reddit.com}{\lr{Reddit} (نظرات دانشجویان را بخوانید)}
\end{itemize}


\subsubsection{بورسیه‌ها و کمک‌های مالی}
علاوه بر بورسیه‌های وابسته به کشور، معمولا دانشگاه‌ها هم برای دانشجویان دارای شرایط خاص بورسیه‌هایی در نظر میگیرند. میزان و تعداد افراد 
مورد پوشش این بورسیه‌ها بستگی به بودجه دانشگاه‌ها دارد و عموما محدود است. برای اطلاع از این بورسیه‌ها می‌توانید از سایت دانشگاه‌ها وارد بخش 
مربوط به هزینه‌ها و بورسیه‌ها شوید و آن‌ها را بررسی کنید. چند مثال از این بخش در سایت دانشگاه‌های مختلف:


\begin{itemize}[leftmargin=1.5cm]
  \item \href{https://www.polimi.it/en/prospective-students/how-much-does-it-cost/scholarships}{\lr{Politecnico di Milano}}
  \item \href{https://www.kuleuven.be/scholarships/year}{\lr{KU Leuven}} 
  \item \href{https://erasmus-plus.ec.europa.eu/opportunities/individuals/students/erasmus-mundus-joint-masters}{\lr{Erasmus Mundus Scholarships}} 
  \item \href{https://www.daad.de/en/studying-in-germany/scholarships/daad-scholarships/}{\lr{DAAD Scholarships}} 
  
\end{itemize}


\section{انواع پروگرم‌ها در پذیرش دادن}
پروگرم‌های مختلف در دانشگاه‌ها روش‌های مختلفی برای انتخاب پذیرفته شدگان از بین کسانی که برای یک پروگرم خاص اپلای کرده‌اند دارند. 
دو دسته‌بندی کلی شامل کمیته‌ای و استاد محور می‌شود که همانطور که از اسم‌شان پیداست در مورد اول انتخاب به عهده‌ کمیته‌ای از اساتید
گذاشته می‌شود و در روش دوم یک استاد خاص که قرار است سرپرست دانشجو باشد نظر اصلی درباره پذیرفته شدن دانشجو را می‌دهد. 

\subsection{کمیته‌ای}
در این  روش برای ارزیابی دانشجویانی که به یک پروگرم اپلای کرده‌اند مدارک آن فرد در کمیته‌ای از اساتید آن دانشکده 
مورد بررسی قرار می‌گیرد و نظر نهایی در یک شورا داده می‌شود. معیار‌های ارزیابی معمولا به طور دقیق ذکر نمی‌شود ولی ترکیبی از انگیزه که در انگیزه‌نامه 
درباره آن نوشته‌اید، شایستگی تحصیلی که می‌تواند شامل نمرات، کار پژوهشی و توصیه‌نامه‌های شما باشد. بهترین منبع برای فهمیدن 
این معیار‌ها سایت همان دانشگاه و پروگرمی است که برای آن اپلای می‌کنید. به عنوان مثال می‌توانید معیار‌های مورد استفاده در پذیرش در 
رشته \lr{Machine Learning, Data Science and Artificial Intelligence} در دانشگاه \lr{Aalto} را مطالعه کنید.

\begin{itemize}[leftmargin=1.5cm]
\item \href{https://www.aalto.fi/en/study-options/machine-learning-data-science-and-artificial-intelligence-master-of-science-technology}{\lr{Machine Learning, Data Science and Artificial Intelligence}}
\end{itemize}

این نوع پروگرم‌ها در کشور‌های اروپایی متداول هستند، به خصوص برای مقطع ارشد که بخش عمده نوشتن و تحقیق پایان‌نامه در سال دوم 
تحصیل انجام می‌شود و دانشجو می‌تواند هنگام تحصیل و در سال اول استاد مورد نظر خود را پیدا کند. تعداد پروگرم‌های کمیته‌محوری که دارای 
بورسیه و فاند می‌باشند کمتر می‌باشد اما این به این معنی نیست که چنین پروگرم‌ها و دانشگاه‌هایی وجود نداند. 

\subsection{استاد محور}
برای پذیرش در این پروگرم‌ها قبل از ارائه مدارک به دانشگاه برای صدور نامه رسمی پذیرش باید از طریق مکاتبه یا روش‌های مشابه، هماهنگی‌های لازم
با یک استاد از آن دانشکده را انجام داده باشید. به عنوان مثال از طریق ایمیل، علاقه خود را به زمینه تحقیقاتی استاد اعلام می‌کنید و استاد 
مورد نظر ممکن است روش‌های مختلفی برای ارزیابی شما داشته باشد، مثل بررسی رزومه، ریزنمرات، مصاحبه‌فنی و موارد مشابه. در نهایت در صورتی که 
استاد طی این ارزیابی‌ها به شما تایید بدهد که می‌توانید درخواست پذیرش خود را به دانشگاه بدهید  مدارک خود را به دانشگاه ارائه خواهید کرد و 
استاد به دانشگاه اعلام خواهد کرد که شما را به عنوان دانشجوی خود انتخاب کرده و دانشگاه با انجام ارزیابی‌های محدود‌تری نسبت به پروگرم‌های کمیته‌محور
نظر خود را اعلام خواهد کرد. معمولا دانشگاه در صورت نظر مساعد استاد، نظر مخالف ارائه نخواهد کرد و می‌توانید پذیرش خود را بگیرید. 
این پروگرم‌ها معمولا دارای فاند هستند زیرا دلیل ارزیابی شدن شما توسط استاد این است که استاد در حال برنامه‌ریزی استفاده از گرنت و بودجه تحقیقاتی
خود برای جذب دانشجو است و به همین دلیل نظر استاد اهمیت بیشتری در فرایند قبولی دارد. این پروگرم‌ها در کانادا، آمریکا و مقاطع دکتری و پست‌داک 
متداول‌تر هستند ولی مثل همیشه استثنا‌هایی هم وجود دارند. 


\section{مراحل گرفتن پذیرش - کمیته‌محور}
در این مرحله فرض شده است که دانشگاه‌ و پروگرم‌های مد نظر خود را پیدا کرده، یادداشت کرده‌اید و از ددلاین‌های آن‌ها با خیر هستید. 
توصیه می‌شود برای جلوگیری از فراموشی و به هم ریختگی یک فایل اکسل یا گوگل‌شیت برای نگه‌داشتن پروگرم‌های مختلف، لینک وبسایت آن‌ها 
ددلاین، اطلاعات بورسیه و موارد مشابه داشته باشید. 

\subsection{آماده‌سازی مدارک}
همانطور که قبلا توضیح دادیم، دانشگاه‌های کمیته‌محور معمولا مجموعه‌ای از مدارک را می‌خواهند که تقریبا بین اکثر آنها ثابت است. این مجموعه
مجموعه معمولا شامل \hyperref[ssec:transcripts]{ریزنمرات}، \hyperref[ssec:cv]{رزومه}، \hyperref[ssec:lang-cert]{مدرک زبان}، \hyperref[ssec:sop]{انگیزه‌نامه} و \hyperref[ssec:recom]{توصیه‌نامه} است. که می‌توانید در \hyperref[sec:admission-docs]{بخش مدارک} درباره آنها و نحوه آماده‌سازیشان مطالعه کنید. 
بیشتر این مدارک باید از قبل آماده شده باشند، ولی برخی مدارک مثل توصیه‌نامه و انگیزه‌نامه ممکن است زمان بیشتری بگیرند یا مثل مدرک زبان نیاز به گرفتن نوبت از قبل یا صرف زمان طولانی برای آموزش داشته باشند. پر کردن توصیه‌نامه 
چون توسط استاد انجام می‌شود ممکن است چند روز طول بکشد، بنابرین در زمان‌بندی خود این مورد را نظر داشته باشید. 

\subsection{ثبت اپلیکیشن}
در کشور‌های مختلف ثبت اپلیکیشن برای یک پروگرم می‌تواند فرایند‌های مختلفی داشته باشد، مثلا بعضی دانشگاه‌ها در فرانسه و آلمان 
از پورتال‌های مرکزی به نام‌های \lr{Uniassist} و \lr{Campus France} استفاده می‌کنند، پس ابتدا بررسی کنید که اپلیکیشن خود را در چه وبسایتی 
باید ثبت کنید. 
پس از مشخص شدن این مورد، مدارک و مشخصات مورد نیاز را در وبسایت مربوطه وارد می‌کنید و در صورت نیاز ایمیل اساتید مورد نظرتان برای دریافت 
توصیه‌نامه را ثبت می‌کنید، در ادامه برای ثبت نهایی اپلیکیشن  معمولا باید مبلغی به عنوان \lr{Application Fee} پرداخت کنید. 
این کار را می توانید خودتان در صورت داشتن کارت اعتباری خارجی، دوستان یا آشناهای ساکن خارج از ایران و یا توسط موسسات پرداخت خارجی
انجام دهید. موسسات معروف در این زمینه شامل شریف‌ویزا، شریفی‌منش، ایرانیکارت، یزد‌پیمنت و ده‌ها موسسه دیگر هستند. مبلغ مورد نیاز 
برای اپلیکیشن فی دانشگاه‌های اروپایی معمولا بین 80 تا 100 یورو است. 


\subsection{نتیجه}
معمولا نتیجه یا مواردی که نیاز به انجام کاری از طرف شما داشته باشند از طریق ایمیل اطلاع‌رسانی می‌شوند. دانشگاه‌های مختلف 
بین چند هفته تا چند ماه برای ارزیابی داوطلبان زمان نیاز دارند، زمان تخمینی ارزیابی را از وبسایت دانشگاه‌ها بررسی کنید. 
پس از گذشت زمان لازم و اطلاع پیدا کردن از نتیجه، در صورت قبولی معمولا باید قبولی را \lr{Accept} یا \lr{Withdraw/Reject} کنید. که 
در صورت پذیرفتن آن یک نامه \lr{Admission} دریافت خواهید کرد که در مراحل بعدی مثل دریافت ویزا و معافیت تحصیلی خارج به آن 
نیاز خواهید داشت. 


\section{مراحل گرفتن پذیرش - استادمحور}
همانطور که توضیح دادیم برای پذیرش در بعضی پروگرم‌ها ابتدا باید استادی پیدا کنید که حاضر باشد به عنوان \lr{supervisor} شما در 
مقطع مورد نظر فعالیت کند. بر خلاف پروگرم‌های کمیته‌محور عامل اصلی ارزیابی شما \textbf{در ابتدا} یک استاد خاص و نه کمیته یک دانشگاه است 
و به همین دلیل اینکه چه مدارکی نیاز دارید به استاد بستگی دارد ولی اغلب اساتید به رزومه و ریزنمرات (ترجمه شده) اکتفا می کنند. 

\subsection{پیدا کردن اساتید}
برای پیدا کردن سوپروایزر ابتدا باید به لیستی از اساتید دسترسی پیدا کنید و زمینه پژوهشی هر یک را به طور کلی مورد بررسی قرار دهید، 
این کار را می‌توانید از منابع مختلفی انجام بدهید. 

\begin{enumerate}
  \item سایت دانشکده‌ هر دانشگاه (مثال:‌ \href{https://www.ualberta.ca/en/computing-science/faculty-and-staff/facultyold.html}{\lr{Alberta CS Faculty}})
  \item \href{https://csrankings.com}{\lr{CS Rankings}}
  \item \href{https://airankings.com}{\lr{AI Rankings}}
  \item لیست آماده اساتید که از سال‌های قبل و توسط دیگر دانشجویان تهیه شده
\end{enumerate}

\begin{tipbox}
فقط به موارد ۲، ۳ و ۴ اکتفا نکنید، به‌روز ترین اطلاعات و مرجع اصلی سایت دانشکده مورد نظر شماست. جابجایی اساتید، بازنشسته شدن 
یا اضافه شدن استاد‌های جدید بسیار رایج است و برخی منابع ممکن است دیر به دیر بروزرسانی شوند.
\end{tipbox}

با دسترسی به لیست‌های اساتید می‌توانید بسته به زمینه‌پژوهشی آن‌ها، وبسایت و \lr{homepage} هر استاد تصمیم بگیرید که چقدر 
به کار کردن با آن استاد اشتیاق دارید. در ضمن بعضی از اساتید در وبسایت خودشان یا آزمایشگاهشان بخشی تحت عنوان \lr{Open Positions/Jobs/Prospective Students} 
قرار می‌دهند که در آن اطلاعات موقعیت‌های باز و خواستن یا نخواستن دانشجو را قرار می‌دهند. در همین پروسه می‌توانید 
اساتیدی که به زمینه پژوهشی شما نزدیک هستند را مشخص کنید و در یک لیست دوم وارد کنید یا به نحوی علامت گذاری کنید و در همین حین می‌توانید مرحله بعدی
را به صورت موازی با این مرحله شروع نمایید.


\subsection{مکاتبه}
وبسایت آزمایشگاه و یا \lr{homepage} استاد مورد نظر را با دقت مطالعه کنید. بعضی اساتید توضیحات خاصی برای 
دانشجویانی که می‌خواهند با آنها کار کنند در وبسایتشان قرار می‌دهند مثلا شاید از شما بخواهند به جای ایمیل زدن یک \lr{Google Form} را 
پر کنید، یا کلمه خاصی را در \lr{subject} ایمیل خود قرار بدهید (جلوگیری از اسپم) یا مدرک خاصی را در ایمیل خود داشته باشید. 

پس از مطالعه و مطمئن شدن از این موارد، در صورتی که میخواهید با استاد مکاتبه کنید باید به استاد یک ایمیل بنویسید و به صورت 
بسیار کوتاه توضیح بدهید که \textbf{چرا سوابق پژوهشی/کاری و علایق شما به استاد بسیار نزدیک است و شما انگیزه ‌زیادی برای کار با او دارید}. 
هیچ فرمت یا تمپلیت خاصی برای ایمیل زدن به اساتید وجود ندارد که همه جا قابل استفاده باشد، همچنین استفاده از تمپلیت توصیه نمی‌شود و 
ایمیل شما در صورت تمپلیتی بودن با احتمال بیشتری نادیده گرفته خواهد شد. 

مثالی از یک ایمیل که به گرفتن مصاحبه منجر شد:

\begin{framed}
\begin{latin}

\raggedleft               

Dear Professor [Last Name],

I'm a final-year BSc student in Computer Engineering at [University], ranked [Rank] in [Country], with a GPA of [GPA]/4. I am writing to express my interest in working in the [Lab] as a master's student.

I found your research on medical imaging fascinating, especially your work titled "[Paper Title]". In my recent paper ([paper URL]), I used a similar method of utilizing [Method] to [Application]. I have also worked on computer vision foundations in another project, [Project Name], working with vision models, OpenCV, and face detection.

I find the research on medical image analysis influential and very promising, and aligned with my interests and background. I hope to discuss the possibility of joining your lab over a short online meeting. Please find my CV and transcripts attached for your consideration.

Many thanks,

[Name]
\end{latin}
\end{framed}

برای کسب اطلاعات بیشتر درباره نحوه ایمیل زدن به اساتید می‌توانید از منابع زیر استفاده کنید:

\begin{itemize}
  \item \href{https://www.youtube.com/watch?v=xXZyZzW86LY}{EZ Apply part 1}
  \item \href{https://www.youtube.com/watch?v=m6Kr2Bf5zyk}{EZ Apply part 2}
\end{itemize}

\subsection{مصاحبه/تسک}
معمولا در صورتی که استادی با خواندن ایمیل شما تصمیم به بررسیتان به عنوان یک گزینه اولیه بگیرد، یک مرحله بیشتر ارزیابی انجام خواهد داد. 
این ارزیابی ممکن است به شکل انجام یک تسک مثلا خلاصه کردن یکی از مقالات استاد، پیاده‌سازی یک پروژه ساده، یا انجام یک مصاحبه فنی باشد. 

نوع، شکل و سختی این مصاحبه یا تسک‌ها کاملا به استاد بستگی دارد و چندان قابل پیش‌بینی نیست با این حال سوال کردن از استاد 
درباره خود مصاحبه و اینکه باید چیز خاصی برای آن آماده کنید یا نه می‌تواند به شما کمک کند. 

قبل از مصاحبه سوال‌هایی که درباره پروگرم مقصد دارید مثل شرایط فاندینگ، پروژه‌هایی که استاد روی آن کار می کند و موارد این چنینی را یادداشت کنید
تا در فرصت مناسب در مصاحبه یا بعد از آن از استاد بپرسید

در این مرحله خروجی‌های مختلفی ممکن است اتفاق بیوفتد، برخی اساتید مراحل بیشتری در ازریابی به کار می‌برند، بعضی ممکن است به شما تایید 
نهایی را بدهند و بگویند در صورت اپلای به دانشگاه به آن‌ها اطلاع دهید تا پشتیبانی از شما را به دانشگاه اعلام کنند و بعضی هم جواب مبهم‌تری به شما
خواهند داد. 

\subsection{اپلای به دانشگاه}
پس از گرفتن جواب استاد در این مرحله لااقل برای مقطع ارشد معمولا باید اطلاعات خود را در سایت دانشگاه هم ثبت کنید تا دانشگاه برخی 
ارزیابی‌های اولیه را انجام دهد. در دانشگاه‌هایی که پروگرم‌های استاد محور دارند این بخش از اهمیت کمتری برخوردار است و در صورت تایید دادن 
استاد احتمال اینکه دانشگاه شما را رد کند کمتر است، با این حال در ثبت کردن مدارک دقت کافی را به کار ببرید. 
از اآنجایی که این مرحله بسیار شبیه اپلای به دانشگاه‌های کمیته‌محور است از توضیح اضافه درباره آن خودداری می‌کنیم. 



\section{مدارک}\label{sec:admission-docs}

\subsection{رزومه (\lr{CV})}\label{ssec:cv}
رزومه یکی از مهم‌ترین موارد در ارزیابی شما است. توصیه می‌شود که برای فرصت‌های کاری و تحصیلی نسخه‌های متفاوتی از رزمه‌خود را استفاده کنید زیرا در هر یک
از این موقعیت‌ها نهاد ارزیابی کننده معیار‌های متفاوتی برای ارزیابی دارد، مثلا در رزومه‌های کاری تجربه‌های کاری را اول قرار می‌دهید و در رزومه تحصیلی دانشگاه محل
تحصیل شما، تجربیات کار کردن با اساتید و موارد مشابه می تواند اهمیت بیشتری از تجربه کاری داشته باشد. 
همچنین بهتر است که رزومه‌ و مدارکی که در زمان درخواست پذیرش برای یک دانشگاه استفاده می‌کنید را در یک پوشه مختص آن دانشگاه نگه‌داری کنید. 
برای اپلای در مقطع ارشد بهتر است که رزومه از ۲ صفحه و برای دکتری از ۳ الی ۴ صفحه بیشتر نشود. 

\begin{tipbox}
  کمالگرا نباشید نسخه اول تقریبا همیشه بد خواهد بود. به مرور آن را بهتر خواهید کرد.
\end{tipbox}


در اولین فرصتی که می‌توانید یک نسخه اولیه از رزومه خود در یک غالب مناسب تهیه کنید و سعی کنید به مرور زمان با به اشتراک گذاشتن آن بین دوستان و سال‌بالایی‌ها
کیفیت آن را بالا ببرید. استفاده از \LaTeX برای نوشتن رزومه توصیه می‌شود. می‌توانید از قالب‌های آماده استفاده کنید، لیست زیر تعدادی
از قالب‌‌های مناسب را ارائه می‌دهد:

\begin{itemize}[leftmargin=1.5cm]
\item \href{https://www.overleaf.com/latex/templates/harshibars-resume/sbcyynmtpnyd}{\lr{Template 1}}
\item \href{https://www.overleaf.com/latex/templates/autocv/scfvqfpxncwb}{\lr{Template 2}}
\item \href{https://www.overleaf.com/latex/templates/plantilla-cv-espanol/ttbgvmyjcpfb}{\lr{Template 3}}
\item \href{https://www.overleaf.com/articles/yash-thakurs-curriculum-vitae/grsjzdtqqrzj}{\lr{Template 4}}
\item \href{https://www.latextemplates.com/template/medium-length-graduate-cv}{\lr{Template 5}}
\end{itemize}

چیزهایی که توصیه می‌شود در رزومه ذکر کنید (به ترتیب):

\begin{itemize}[leftmargin=1.5cm]
  \item معدل و دانشگاه
  \item سابقه کار تحقیقاتی و \lr{Research Assistant}
  \item مقاله و کار پژوهشی
  \item مهارت‌های فنی و مرتبط
  \item سابقه کاری مرتبط
  \item پروژه‌ها و لینک‌های مربوط به آن‌ها
  \item تجربه تدریسیاری
  \item نمره زبان یا تاریخی که احتمالا آزمون زبان می دهید
\end{itemize}
توجه کنید که به طور مثال اگر سابقه کاری ندارید یا پروژه‌ای که مرتبط به پروگرمی که برای آن اپلای می‌کنید نیست بهتر است آن را نیاورید،
شلوغ کردن رزومه باعث توجه کمتر مخاطب خواهد شد. یک رزومه خلوت که موار مرتبط در آن آمده بهتر از یک رزومه پر از دستاورد‌های بی‌ربط است. 

برای راهنمایی دقیق‌تر و کامل‌تر می‌توانید از منابع زیر استفاده کنید:

\begin{itemize}[leftmargin=1.5cm]
\item \href{https://github.com/shaily99/advice#cv}{\lr{CV Examples and Advice}}
\item \href{https://www.youtube.com/playlist?list=PL3rXObcfMsILRH6bypjHRs5XS6ANLc2oI}{\lr{EZApply - CV}}
\end{itemize}


\subsection{انگیزه‌نامه (\lr{Statement of Purpose})}\label{ssec:sop}
یکی از مهمترین مدارکی که هنگام ارزیابی شما بررسی می‌شود انگیزه‌نامه است. در انگیزه‌نامه توضیح می‌دهید که \textbf{سابقه و پیشینه‌تان} چه بوده و 
در ادامه باید خواننده را قانع کنید که این سوابق چطور باعث شده که شما شما انگیزه زیادی برای تحصیل در مقطع ارشد/دکتری در این 
رشته، دانشگاه و پروگرم خاص داشته باشید و \textbf{کاملا برای آن مناسب} باشید. هرچقدر که انگیزه‌نامه بتواند گذشته ‌شما را به تصمیم امروزتان برای 
تحصیل در یک رشته و پروگرم خاص متصل کند و داستان بهتری در این زمینه بگوید برایتان بهتر خواهد بود. نوشتن انگیزه‌نامه نسبت به رزومه زمان بیشتری می‌برد،
بنابرین سعی کنید برای نوشتن، بازخوانی و بهبود نسخه اولیه آن حدود یک الی دو هفته زمان داشته باشید. پس از این زمان نوشتن نسخه‌های آینده 
زمان کمتری خواهند برد. 

انگیزه‌نامه را می‌توان به شیوه‌های متفاوتی نوشت ولی به طور کلی یک انگیزه نامه خوب باید شامل موارد زیر باشد:

\subsubsection{مقدمه کوتاه \lr{(Introduction)}}
یک توضیح کوتاه در حد یک پاراگراف از اینکه چرا به رشته و زمینه‌ای که برای آن اپلای می‌کنید علاقه‌مند هستید و می‌خواهید آن را در مقطع بالاتر 
ادامه بدهید. 

\subsubsection{پیشینه \lr{(Background)}}
سوابق و تجربیات مرتبط با پروگرم هدف. این بخش می‌تواند شامل مقالات، کار پژوهشی و آزمایشگاهی، پایان‌نامه، تدریسیاری، سوابق کاری، و سایر مواردی باشد که شما را 
از لحاظ آکادمیک بالاتر از سایرین قرار می‌دهد. نکته‌ای که باید در این بخش به خاطر بسپارید نشان دادن یک ترتیب منطقی بین تجربیاتتان و آوردن تجربیات مرتبط 
و پرهیز از اضافه‌گویی است. مثلا اگر برای یک پروگرم مرتبط با شبکه جملات زیر می‌تواند برای انگیزه‌نامه مناسب باشد:

\begin{framed}
\lr{
...Building towards this vision, I focused on my coursework and maintained a GPA of X. I sought to apply my knowledge to novel 
challenges. Therefore, I joined Amirkabir's X Lab for a project on Y. During this 
project, I proposed methods to assess Z. To document our findings, we co-authored a manuscript (cite the manuscript), published at 
VERY GOOD JOURNAL 2026. This trajectory equipped me with the necessary skills for the more complex research required in a master's program.
}
\end{framed}


\subsubsection{اهداف \lr{(Goals)}}
در این بخش باید هدف و انگیزه خود را از تحصیل در \textbf{مقطعی که برای آن اپلای می‌کنید} را به طور خاص برای
این \textbf{رشته} و این \textbf{دانشگاه} توضیح دهید. هدف شما باید دقیق، قانع‌کننده و منطقی باشد و باید از نوشتن
جملات عمومی پرهیز کنید. 

\begin{warningbox}
  جملات زیر مثال‌هایی از جملات عمومی (\textbf{generic}) هستند که هیچ حقیقت خاصی را درباره شما و دلیل
  علاقه‌مند بودنتان به یک پروگرم خاص بیان نمی‌کنند. این جملات باعث قانع شدن مخاطب و قائل شدن فرق بین شما 
  و یک فرد دیگر نخواهند شد. 

  \begin{itemize}
    \item  \lr{I realized that I like to follow Aritficial Intelligence for my MSc.}
    \item \lr{I sought to expand my knowledge on cyber security, therefore I decided to apply for a PhD in this field.}
  \end{itemize}
\end{warningbox}

\subsubsection{مطابقت با دانشگاه/استاد \lr{(Alignment/Program Fit/Faculty Fit)}}
بهتر است برای هر پروگرم یک نسخه خاص از \lr{SoP} خود بنویسید، البته نیاز نیست که انگیزه‌نامه کاملا از صفر نوشته شود و می‌توانید یک نسخه اولیه و عمومی از 
انگیزه‌نامه خود داشته باشید و آن را برای پروگرم‌های مختلف مناسب‌سازی کنید. 

\subsubsection{پروگرم‌های کمیته‌محور}
در این پروگرم‌ها تاکید این بخش روی مزیت‌های دانشگاه، امکانات آن، رتبه دانشکده به خصوصی و سایر مواردی است که آن پروگرم را به \textbf{طور خاص برای شما}
بسیار مناسب می‌کند. علاوه بر امکانات نکته مهمی که می‌توانید ذکر کنید (مخصوصا در مقطع ارشد) کورس‌های یک پروگرم است که در راستای رسیدن 
به اهداف تحصیلی و کاریتان به شما کمک می‌کند. به عنوان مثال:

\begin{framed}
\lr{
...Although I took Operating System Concepts course during my BSc and also served as a TA, my knowledge on memory bandwidth optimization is
shallow. The Advanced Operating Systems Course (COMP 7321) fills this exact knowledge gap, which will later serve me to conduct my MSc-level 
research on addressing the memory bottleneck when serving LLMs.
}
\end{framed}

این متن توضیح می‌دهد که شما با اینکه آشنایی اولیه با یک مبحث را دارید، درس‌هایی که در این پروگرم خاص ارائه می‌شود به شما کمک می‌کند که 
بخش‌هایی از دانشتان که برای انجام پژوهش مورد نظرتان در مقطع ارشد کافی نیست را پوشش دهید. 


\subsubsection{پروگرم‌های استاد محور}
در این پروگرم‌ها می‌توانید مثل قبل به درس‌هایی که ارائه می‌شود اشاره کنید ولی موردی که اهمیت بیشتری دارد ذکر کردن نام استاد(های) مد نظرتان 
در آن دانشگاه به عنوان استاد راهنما است. در صورتی که از قبل با یک استاد خاص هماهنگ کرده اید و تایید او را دریافت کرده‌اید در این بخش 
به \textbf{نام، پژوهش‌ها و حتی یکی از کارهای پژوهشی } استاد اشاره کنید و توضیح دهید که چرا این استاد کاملا با اهداف شما هماهنگ است. 
در صورتی که با استاد خاصی به توافق نرسیده‌اید یا می‌خواهید به صورت \textbf{blind} برای یک پروگرم استاد محور اپلای کنید، می‌توانید همین کار
را برای دو یا نهایتا سه استاد انجام بدهید و توضیح بدهید که چرا پژوهش‌های آنها با علایق و سلیقه شما کاملا سازگار است، هرچقدر بتوانید توضیح دقیق‌تری
در این بخش بدهید بهتر است، مثلا یک مقاله یا فعالیت پژوهشی خودتان را ذکر کنید که با کارهای پژوهشی استاد همسو است یا مثلا توضیح دهید که 
سوابق شغلی و واحد‌هایی که پاس کرده‌اید (بهتر است اسم ذکر کنید) با فعالیت‌های استاد همسو است. 

منابع مناسب برای مطالعه و یادگیری نوشتن انگیزه‌نامه:

\begin{itemize}[leftmargin=1.5cm]
\item \href{https://www.reddit.com/r/StatementOfPurpose/comments/9k40n9/sop_tips_thread/}{\lr{Reddit}}
\item \href{https://www.reddit.com/r/gradadmissions/comments/zn0u1d/resources_for_writing_a_statement_of_purpose/}{\lr{SOP Resources (Cornell, Berkeley, Northeastern, UCSB, …)}}
\item \href{https://cs-sop.notion.site/cs-sop/CS-PhD-Statements-of-Purpose-df39955313834889b7ac5411c37b958d}{\lr{SOP Bank}}
\end{itemize}

\subsection{توصیه‌نامه (\lr{Recommendation Letter})}\label{ssec:recom}
همانطور که از اسم آن مشخص است، توصیه‌نامه معمولا شامل یک متن از پیش نوشته شده درباره شما توسط استاد و گاهی پر کردن یک پرسشنامه 
درباره توانایی‌های مختلف شما می‌شود. معمولا توصیه‌نامه باید به دست خود استاد و در سایت دانشگاه مقصد آپلود شده یا توسط استاد برای آنجا ایمیل
شود ولی مواردی هم وجود دارد که اجازه می‌دهد \lr{PDF} توصیه‌نامه توسط خود دانشجو در هنگام پر کردن اپلیکیشن ثبت شود. 
به طور کلی نوشتن توصیه‌نامه به عهده استاد است ولی برخی اساتید از دانشجو یک نسخه \lr{draft} برای توصیه‌نامه می‌خواهند که بعدا خودشان 
در آن تغییراتی ایجاد می‌کنند. 

پروگرم‌ها بسته به نوع بین ۰ تا ۳ توصیه‌نامه نیاز دارند. معمولا پروگرم‌های کانادایی ۳ یا ۲ توصیه‌نامه، اروپایی ۲ یا یک توصیه‌نامه نیاز دارند و 
برخی پروگرم‌ها هم توصیه‌نامه را به عنوان موردی اختیاری ذکر می‌کنند. 

\begin{tipbox}
  \begin{itemize}
    \item   هنگام شروع فرایند اپلای یک یادداشت یا شیت به خصوصی برای توصیه‌نامه‌هایی که از اساتید گرفتید و جایی که آن‌ها را ثبت کرده‌اند بسازید
  در ادامه فرایند این یادداشت به شما کمک می‌کند که توصیه‌نامه‌های خود را بهتر مدیریت کنید.
    \item در همه انواع توصیه‌نامه‌ها بهتر است شخص نویسنده اشاره کند که از چه زمان شما را می‌شناسد (مثلا یک سال دو سال یا بیشتر) ذکر کردن
    این مورد کمک می‌کند که کسی که توصیه‌نامه را بررسی می‌کند تخمین بهتری از میزان شناخت نویسنده از شما داشته باشد
    \item اگر یک پروگرم بیشتر از یک توصیه‌نامه نیاز دارد، سعی کنید در صورتی که می‌توانید که از بیشتر از یک نوع از توصیه‌نامه‌های توضیح داده شده استفاده کنید، استفاده 
    از دو توصیه‌نامه از یک نوع(مثلا تدریسیاری) اطلاعات اضافه‌ای به فرایند بررسی پرونده شما اضافه نمی‌کند. 
  \end{itemize}

\end{tipbox}

در ادامه سه نوع مهم توصیه‌نامه را بررسی می‌کنیم:

\subsubsection{توصیه‌نامه ریسرچی}
این نوع توصیه‌نامه باید توسط استاد/شخص/سازمانی نوشته شود که با او کار پژوهشی انجام داده‌اید. این کارهای پژوهشی می‌تواند شامل 
نوشتن مقاله، تحقیق روی یک زمینه خاص، انجام پروژه کارشناسی یا پایان‌نامه ارشد و موارد مشابه باشد. متن این نامه باید شامل جزئیاتی از 
زمینه انجام پژوهش، نقش شما در انجام آن و توانایی‌هایی باشد که در انجام پروژه و از خود نشان داده‌اید مثلا مهارت بالا در حل مسئله
یا مهارت در نوشتن متن‌های علمی و موارد مشابه.

\subsubsection{توصیه‌نامه تدریسیاری}
بسیاری از استاد‌ها به دانشجویانی که در یکی از درس‌های استاد تدریسیار آن‌ها بوده‌اند توصیه‌نامه می‌دهند، متن این توصیه‌نامه‌ها می‌تواند شامل 
دانش و توانایی شما در آن درس خاص و مهارت‌های نرم مثل تدریس، وظیفه‌شناسی، مهارت ارائه، مهارت رهبری و موارد مشابه می‌باشد. 

\subsubsection{توصیه‌نامه کاری}
\begin{warningbox}
توجه کنید که پیش از استفاده از این نوع توصیه‌نامه باید مطمئن شوید که دانشگاه مقصد این نوع توصیه‌نامه را قبول دارد زیر برخی از پروگرم‌ها
صرفا توصیه‌نامه‌های آکادمیک را از شما قبول میکنند. 
\end{warningbox}

در صورتی که کاری نسبتا مرتبط با پروگرم هدف انجام داده‌اید و در این زمینه سوابق شغلی دارید، ارائه نامه اشتغال به کار که توسط کارفرما نوشته شده
یا پر کردن توصیه‌نامه توسط شخصی در محل کار (مثلا مدیر تیم، مدیر شرکت یا مدیر منابع انسانی) می‌تواند اپلیکیشن شما را تقویت کند. 
در این توصیه‌نامه‌ها به وظایف شغلی، وظیفه‌شناسی و حسن انجام کار، زمانی که در آن شرکت مشغول به کار بوده‌اید و در صورت ربط داشتن به 
موضوع کلی کاری که در آن شرکت انجام می‌دادید اشاره کنید. 


\subsection{ریزنمرات (\lr{Transcripts})}\label{ssec:transcripts}
اکثر پروگرم‌ها موقع اپلای از شما ریزنمرات رسمی نمی‌خواهند و یک نسخه غیررسمی ریزنمرات نیز برای آن‌ها کفایت می‌کند، با این حال بهتر است که ریزنمرات
غیررسمی هم توسط دانشگاه صادر شده باشد و شامل سربرگ یا لوگو دانشگاه باشد. معمولا با مراجعه به آموزش دانشکده خود می‌توانید یک نسخه 
از کارنامه غیر رسمی خود را صرفا با سربرگ دانشگاه تهیه کنید. در ادامه می‌توانید این ریزنمره را ترجمه کرده و از اسکن آن استفاده کنید. 

توصیه می‌شود که در صورت داشتن شرایط مالی با مراجعه به آموزش کل دانشگاه لغو تعهد آموزش رایگان خود را تا ترمی که گذرانده‌اید انجام بدهید
در صورت انجام این کار می‌توانید یک نسخه از ریزنمره رسمی خود که دارای مهر و امضای دانشگاه است نیز تهیه کنید. ولی همانطور که توضیح دادیم 
هنگام ایمیل زدن به اساتید یا اپلای برای پروگرم‌ها معمولا نیازی به ریزنمره رسمی ندارید. 


\subsection{مدرک زبان}\label{ssec:lang-cert}
برخی از پروگرم‌ها موقع ثبت اپلیکیشن، ثبت کردن مدرک زبان را هم الزامی می‌کنند ولی بعضی دیگر فرصت جدا و بیشتری برای ثبت آن می‌دهند. 
توصیه می‌شود که مدرک زبان خود را یک سال قبل از شروع پروگرم مقصد بگیرید، به عنوان مثال اگر برای پروگرم‌هایی اپلای می‌کنید که 
در سپتامبر ۲۰۲۶ شروع می‌شوند، بهتر است قبل از سپتامبر ۲۰۲۵ مدرک زبان خود را بگیرید. 

\subsubsection{IELTS}
آزمونی که معمولا برای پروگرم‌های کارشناسی ارشد و دکتری نیاز دارید \lr{IELTS Academic} است که کمی از \lr{IELTS General} سخت‌تر است. 
نمره این آزمون از ۹ می‌باشد و در زمان نوشتن این متن (۱۴ تیر ۱۴۰۵) این آزمون در ایران برگزار نمی‌شود.
این آزمون در اکثریت قریب به اتفاق دانشگاه‌ها و کشورها معتبر حساب می‌شود\footnote{مثال نقض این موضوع ترکیه است}. نمره‌آی که باید 
در این آزمون بگیرد بستگی به دانشگاه و کشور مقصد دارد ولی جدول \ref{tab:ielts_requirements} می‌تواند کمک کننده باشد:

\begin{table}[h]
\centering
\caption{حداقل نمره‌ی معمول آزمون \lr{IELTS Academic}}
\label{tab:ielts_requirements}
\begin{tabular}{lll}
\toprule
\textbf{نمره} & \textbf{وضعیت} & \textbf{نمونه دانشگاه‌ها} \\
\midrule
\lr{6.0} & حداقل برای برخی دانشگاه‌ها & \lr{Manchester، Bristol} \\
\lr{6.5} & رایج‌ترین حداقل نمره & \lr{UCL، Imperial، Edinburgh، Toronto، Melbourne} \\
\lr{7.0} & دانشگاه‌ها یا رشته‌های رقابتی & \lr{Oxford، LSE، ETH Zurich} (برخی رشته‌ها) \\
\lr{7.5} & رشته‌ها و دانشگاه‌های بسیار رقابتی & \lr{Cambridge}، برخی رشته‌های \lr{Oxford} \\
\bottomrule
\end{tabular}
\end{table}

\begin{tipbox}
در اکثر دانشگاه‌ها، کسب نمره‌ای بالاتر از حداقل موردنیاز تأثیر مستقیمی بر شانس پذیرش ندارد و مدرک زبان صرفاً یک شرط لازم محسوب می‌شود. بنابراین، پس از رسیدن به حداقل نمره، معمولاً بهتر است روی سایر بخش‌های اپلیکیشن مانند رزومه، مقالات و انگیزه‌نامه تمرکز کنید. البته برخی دانشگاه‌ها علاوه بر نمره‌ی کل، برای هر مهارت نیز حداقل نمره تعیین می‌کنند.
\end{tipbox}

\subsubsection{TOEFL}
تکمیل شود